\chapter*{Conclusions of Part~II}
\addcontentsline{toc}{chapter}{Conclusions of Part~II}
\label{ch:part2-concl}

The four chapters in this Part have explored complementary strategies for constraining the dark sector, from connecting DM production mechanisms to their experimental signatures to developing principled inference frameworks for interpreting cosmological data.

On the DM side, Chapter~\ref{ch:minimal-dm-freezein} demonstrated that the freeze-in production mechanism yields a qualitatively different phenomenology when the reheating temperature lies below the DM mass. The exponential suppression of DM production in this regime demands larger portal couplings to match the observed relic abundance, enhancing the DM--electron scattering cross section by a factor that scales as $\sqrt{x_\mathrm{rh}}\,e^{x_\mathrm{rh}}$ with $x_\mathrm{rh} = m_\chi/T_\mathrm{rh}$. This result reframes the freeze-in benchmark as an extended region in parameter space rather than a single curve, with a large portion already being probed by current direct detection experiments. 

While Chapter~\ref{ch:minimal-dm-freezein} showed how production dynamics can enhance laboratory sensitivity, Chapters~\ref{ch:cmb-leptophilic-dm} and~\ref{ch:pbh-memory-burden} approach the problem from the other direction, showcasing the power of cosmological data in constraining DM candidates that remain difficult to probe through terrestrial searches alone. 
I showed that for sub-MeV leptophilic DM the dominant channel for cosmological constraints is not Compton-like scattering but rather the loop-induced decay into two photons, a process often overlooked due to its suppression. Using a cosmologically consistent treatment of the resulting energy deposition and the most recent measurements from the \textit{Planck} satellite, I derived CMB bounds that surpass existing terrestrial constraints by up to two orders of magnitude for pseudo-scalar DM lighter than $\sim 5\,$keV. These two chapters together highlight the complementarity between cosmological probes and laboratory experiments in mapping the viable parameter space for DM.

In Chapter~\ref{ch:pbh-memory-burden}, I extended this approach to a non-particle DM candidate --- primordial black holes~\cite{Zeldovich:1967lct, Hawking:1971ei, Chapline:1975ojl}. The memory-burden mechanism, rooted in quantum-gravitational considerations on how black holes store and process information, proposes that this information content can suppress Hawking evaporation~\cite{Dvali:2018xpy, Dvali:2018ytn, Dvali:2020wft}, and has attracted considerable attention for its potential to open an entirely new mass window for light primordial black holes ($10^4$--$10^{10}$ g) to constitute all of the DM~\cite{Alexandre:2024nuo, Dvali:2024hsb, Thoss:2024hsr}. By subjecting this scenario to CMB and BBN constraints, I showed that the proposed mass window survives only if the onset of stabilization occurs almost immediately, with the black hole retaining all but a fraction $\lesssim 10^{-10}$  of its initial mass before evaporation is suppressed. Any smooth or gradual transition is ruled out.  More broadly, this analysis illustrates how precision cosmological data can provide critical guidance for theoretical explorations, even at the frontier of quantum gravity, by delineating which modifications to semiclassical black hole physics are observationally viable.

The interplay between theoretical modelling and observational interpretation is equally relevant on the dark energy side. In Chapter~\ref{ch:desi-dark-energy}, I critically examined the recent DESI DR2 evidence for a time-varying equation of state. Starting from well-motivated thawing quintessence models inspired by string theory and axion physics, I constructed theory-informed priors on the phenomenological dark energy equation of state parameters $(w_0, w_a)$ by mapping broad, minimally biased priors on the fundamental model parameters to the space directly constrained by observations. Despite the deliberately conservative nature of these theoretical assumptions, the resulting priors differ substantially from the uniform distributions conventionally adopted in cosmological analyses, which assign significant probability to regions of parameter space that no physically motivated model populates. When applied to the combination of DESI and CMB data, the apparent $\sim 3\sigma$ preference for dynamical dark energy is reduced to $\sim 1.3$--$1.8\sigma$, rendering it consistent with a cosmological constant. This result underscores the importance of carefully assessing prior dependence and parameters choices in cosmological inference, particularly when interpreting signals that could be taken as evidence for physics beyond $\Lambda
$CDM.

Looking ahead, each of the directions explored in this Part stands to benefit from near-term observational progress. Ongoing and next-generation direct detection experiments will probe increasingly deeper into the parameter space of feebly interacting DM, motivating the continued development of creative yet minimal production mechanisms that can be confronted with these improving sensitivities. Future CMB measurements and large-scale structure surveys will further tighten constraints on DM interactions --- both with the Standard Model and within the dark sector itself. Fully exploiting these data will require modeling these interactions at the level of non-linear structure formation, where their imprints on small-scale observables such as the matter power spectrum and halo abundances remain theoretically challenging to characterize with the required precision. On the dark energy side, continued DESI data releases and independent surveys will sharpen the picture, but the lessons of Chapter~\ref{ch:desi-dark-energy} suggest that maximizing their physics reach will also require moving beyond the conventional equation-of-state parameterization --- for instance, by working directly with the dark energy density, which connects more immediately to the distance measurements that these surveys actually constrain~\cite{Wang:2025vfb}, a direction I aim to pursue.
